\documentclass[10pt,a4paper]{amsart}
\usepackage{amsmath,amssymb,amsthm,graphicx}
% Retained minimal preamble for the Stacks Project chapter more-algebra.tex (More on Algebra).
% The chapter file pinned-original-source.tex carries only \input{preamble} on its first line, so no
% preamble text lives inside the chapter itself.  The Stacks Project uses one single shared preamble for
% every chapter; the snapshot stages it identically as sources/stacks-chow/preamble.tex and
% sources/stacks-etale-cohomology/preamble.tex.  The block below is copied mechanically, line for line,
% from sources/stacks-chow/preamble.tex lines 14-18 (the Xy-pic setup) and 154-217 (the native theorem
% declarations together with the mathematical macros).  Nothing was added, removed or redefined; the
% wrapper class and the amsmath/amssymb/amsthm/graphicx base packages are supplied by the assembler.
% Chapter: more-algebra.tex -- More on Algebra.
% For commutative diagrams we use Xy-pic
\usepackage[all]{xy}

% We use 2cell for 2-commutative diagrams.
\xyoption{2cell}

% Theorem environments.
%
\theoremstyle{plain}
\newtheorem{theorem}[subsection]{Theorem}
\newtheorem{proposition}[subsection]{Proposition}
\newtheorem{lemma}[subsection]{Lemma}

\theoremstyle{definition}
\newtheorem{definition}[subsection]{Definition}
\newtheorem{example}[subsection]{Example}
\newtheorem{exercise}[subsection]{Exercise}
\newtheorem{situation}[subsection]{Situation}

\theoremstyle{remark}
\newtheorem{remark}[subsection]{Remark}
\newtheorem{remarks}[subsection]{Remarks}

\numberwithin{equation}{subsection}

% Macros
%
\def\lim{\mathop{\mathrm{lim}}\nolimits}
\def\colim{\mathop{\mathrm{colim}}\nolimits}
\def\Spec{\mathop{\mathrm{Spec}}}
\def\Hom{\mathop{\mathrm{Hom}}\nolimits}
\def\Ext{\mathop{\mathrm{Ext}}\nolimits}
\def\SheafHom{\mathop{\mathcal{H}\!\mathit{om}}\nolimits}
\def\SheafExt{\mathop{\mathcal{E}\!\mathit{xt}}\nolimits}
\def\Sch{\mathit{Sch}}
\def\Mor{\mathop{\mathrm{Mor}}\nolimits}
\def\Ob{\mathop{\mathrm{Ob}}\nolimits}
\def\Sh{\mathop{\mathit{Sh}}\nolimits}
\def\NL{\mathop{N\!L}\nolimits}
\def\CH{\mathop{\mathrm{CH}}\nolimits}
\def\proetale{{pro\text{-}\acute{e}tale}}
\def\etale{{\acute{e}tale}}
\def\QCoh{\mathit{QCoh}}
\def\Ker{\mathop{\mathrm{Ker}}}
\def\Im{\mathop{\mathrm{Im}}}
\def\Coker{\mathop{\mathrm{Coker}}}
\def\Coim{\mathop{\mathrm{Coim}}}

% Boxtimes
%
\DeclareMathSymbol{\boxtimes}{\mathbin}{AMSa}{"02}

%
% Macros for moduli stacks/spaces
%
\def\QCohstack{\mathcal{QC}\!\mathit{oh}}
\def\Cohstack{\mathcal{C}\!\mathit{oh}}
\def\Spacesstack{\mathcal{S}\!\mathit{paces}}
\def\Quotfunctor{\mathrm{Quot}}
\def\Hilbfunctor{\mathrm{Hilb}}
\def\Curvesstack{\mathcal{C}\!\mathit{urves}}
\def\Polarizedstack{\mathcal{P}\!\mathit{olarized}}
\def\Complexesstack{\mathcal{C}\!\mathit{omplexes}}
% \Pic is the operator that assigns to X its picard group, usage \Pic(X)
% \Picardstack_{X/B} denotes the Picard stack of X over B
% \Picardfunctor_{X/B} denotes the Picard functor of X over B
\def\Pic{\mathop{\mathrm{Pic}}\nolimits}
\def\Picardstack{\mathcal{P}\!\mathit{ic}}
\def\Picardfunctor{\mathrm{Pic}}
\def\Deformationcategory{\mathcal{D}\!\mathit{ef}}
\usepackage{cleveref}
\crefname{originallemma}{original lemma}{original lemmas}
\Crefname{originallemma}{Original lemma}{Original lemmas}
\makeatletter
\newlabel{more-algebra-lemma-filtered-colimits-henselian}{{0FWT}{source}{}{sourcepointer-1}{}}
\newlabel{more-algebra-lemma-filtered-colimits-henselian@cref}{{[originallemma][0][]0FWT}{[1][0][]source}{}{}{}}
\makeatother
\crefname{originallemma}{original lemma}{original lemmas}
\Crefname{originallemma}{Original lemma}{Original lemmas}
\makeatletter
\newlabel{algebra-lemma-contained-in-radical}{{0AME}{source}{}{sourcepointer-2}{}}
\newlabel{algebra-lemma-contained-in-radical@cref}{{[originallemma][0][]0AME}{[1][0][]source}{}{}{}}
\makeatother
\title{Henselian pairs are preserved under filtered colimits}
\author{The Stacks Project Authors}
\date{stacks-project:more-algebra.tex:local-snapshot-20261012; GNU FDL 1.2 or later}
\begin{document}
\maketitle
\setcounter{subsection}{1}
\noindent\textit{Editorial note.} All delivered mathematical text is the author's own, copied from the pinned Stacks Project chapter file more-algebra.tex; no statement or proof step was rewritten, reordered or removed.\par\medskip
\noindent\textit{Editorial note.} The chapter file pinned-original-source.tex carries only \textbackslash{}input\{preamble\} on its first line, so no preamble text lives inside the chapter. The retained file sources/stacks-more-algebra/retained-preamble.tex is a mechanical, line-for-line copy of the project's single shared preamble as staged in the snapshot (Xy-pic setup plus theorem declarations and mathematical macros). Nothing was added, removed or redefined, and the excerpt serves only as preamble, never as a span coordinate.\par\medskip
\noindent\textit{Editorial note.} Both counted spans start at the \textbackslash{}begin line of their environment, so the statement and the author's proof are delivered whole; surrounding prose of the section is deliberately outside the counted spans and is carried as standard prerequisites instead.\par\medskip
\noindent\textit{Editorial note.} The citation fields (source\_id, source\_version, authors, source\_url, license.name, source\_title) are identical to those of every other card this lane has produced from this chapter, so the batch selector treats them as one source.\par\medskip
\noindent\textit{Allowed standard prerequisites (editorial).}\begin{itemize}
\item Henselian pair, defined through the lifting of factorisations of monic polynomials modulo the ideal. definition used in the same section, source lines 2860-2883, immediately before the counted statement. Supplies the property that has to be preserved and the form in which the lifting is verified.
\item An ideal lies in the Jacobson radical if and only if every element of 1 + I is a unit. the author cites Algebra, Lemma algebra-lemma-contained-in-radical (tag 0AME) at proof line 2897. Turns the invertibility statement proved for 1 + I into the radical condition.
\item Filtered colimits: an element of a filtered colimit is the image of an element of some stage, and the same for a finite system of elements and relations. used at proof lines 2893-2910 to descend f, the factorisation and the Bezout relation to one index j. Is the mechanism that reduces the verification to a single henselian pair.
\end{itemize}\medskip
\ifdefined\hypertarget\hypertarget{sourcepointer-1}{}\fi
\noindent\textit{Source reference.} Original lemma 0FWT: focused-seeds/AGP/sources/stacks-snapshot-20261011/tags.txt, source lines 19858--19858. The lemma's own canonical tag, quoted so that readers can resolve the label inside the delivered text; it points at the counted statement itself and carries no proof content.\par\medskip
\ifdefined\hypertarget\hypertarget{sourcepointer-2}{}\fi
\noindent\textit{Source reference.} Original lemma 0AME: focused-seeds/AGP/sources/stacks-snapshot-20261011/tags.txt, source lines 13402--13402. Navigation pointer to a statement the counted proof refers to but does not contain; nothing load-bearing is delegated to it.\par\medskip
\begin{lemma}
\label{lemma-filtered-colimits-henselian}
The property of being Henselian is preserved under filtered colimits of pairs.
More precisely, let $J$ be a directed set and let $(A_j, I_j)$
be a system of henselian pairs over $J$.
Then $A = \colim A_j$ equipped with the ideal $I = \colim I_j$
is a henselian pair $(A, I)$.
\end{lemma}

\begin{proof}
If $u \in 1 + I$ then for some $j \in J$ we see that $u$ is the
image of some $u_j \in 1 + I_j$. Then $u_j$ is invertible in $A_j$ by
Algebra, Lemma \ref{algebra-lemma-contained-in-radical}
and the assumption that $I_j$ is contained in the Jacobson radical of $A_j$.
Hence $u$ is invertible in $A$. Thus $I$
is contained in the Jacobson radical of $A$ (by the lemma).

\medskip\noindent
Let $f \in A[T]$ be a monic polynomial and let $\overline{f} = g_0 h_0$
be a factorization with $g_0, h_0 \in A/I[T]$ monic generating the unit
ideal in $A/I[T]$. Write $1 = g_0 g'_0 + h_0 h'_0$ for some
$g'_0, h'_0 \in A/I[T]$. Since $A = \colim A_j$ and $A/I = \colim A_j/I_j$
are filtered colimits we can find a $j \in J$ and $f_j \in A_j$ and
a factorization $\overline{f}_j = g_{j, 0} h_{j, 0}$
with $g_{j, 0}, h_{j, 0} \in A_j/I_j[T]$ monic
and $1 = g_{j, 0} g'_{j, 0} + h_{j, 0} h'_{j, 0}$
for some $g'_{j, 0}, h'_{j, 0} \in A_j/I_j[T]$ with
$f_j, g_{j, 0}, h_{j, 0}, g'_{j, 0}, h'_{j, 0}$
mapping to $f, g_0, h_0, g'_0, h'_0$.
Since $(A_j, I_j)$ is a henselian pair, we can lift
$\overline{f}_j = g_{j, 0} h_{j, 0}$ to a factorization
over $A_j$ and taking the image in $A$ we obtain a
corresponding factorization in $A$. Hence $(A, I)$ is henselian.
\end{proof}

\end{document}
